\section{总结与延伸}

\subsection{全课知识图谱}

本节课建立了从深度学习基础到卷积神经网络的完整知识链：

\begin{center}
\begin{tikzpicture}[node distance=0.8cm, every node/.style={draw, rounded corners, minimum width=3.5cm, minimum height=0.7cm, font=\small}]
\node (review) {深度学习基础回顾};
\node (feature) [below=of review] {特征工程 vs 端到端学习};
\node (conv) [below=of feature] {卷积层与池化层};
\node (arch) [below=of conv] {CNN 架构设计};
\node (alex) [below=of arch] {AlexNet 与深度学习革命};
\draw[->, thick] (review) -- (feature);
\draw[->, thick] (feature) -- (conv);
\draw[->, thick] (conv) -- (arch);
\draw[->, thick] (arch) -- (alex);
\node[draw=none, right=0.8cm of review, text width=5.5cm, font=\scriptsize] {线性分类器、损失函数、优化、计算图、反向传播};
\node[draw=none, right=0.8cm of feature, text width=5.5cm, font=\scriptsize] {颜色直方图、HoG、端到端学习优势};
\node[draw=none, right=0.8cm of conv, text width=5.5cm, font=\scriptsize] {局部连接、权值共享、感受野、步幅、填充};
\node[draw=none, right=0.8cm of arch, text width=5.5cm, font=\scriptsize] {Conv-ReLU-Pool 模式、通道数递增、空间递减};
\node[draw=none, right=0.8cm of alex, text width=5.5cm, font=\scriptsize] {GPU 训练、ImageNet、CNN vs Transformer};
\end{tikzpicture}
\end{center}

\subsection{关键 Takeaways}

\begin{importantbox}{五条核心原则}
\begin{enumerate}
    \item \textbf{端到端学习优于手工特征}：让梯度下降从数据中学习特征表示，而非依赖人类直觉
    \item \textbf{卷积利用了图像的先验知识}：局部性（每个输出只看一小块区域）和平移不变性（同一滤波器在所有位置共享）
    \item \textbf{层级特征自动涌现}：浅层学边缘和颜色，深层学纹理和物体部件，无需人工设计
    \item \textbf{非线性是深度网络的灵魂}：没有激活函数，再深的网络也等价于一层线性变换
    \item \textbf{设计模式}：空间尺寸逐步缩小 + 通道数逐步增加 = 从像素到语义的特征金字塔
\end{enumerate}
\end{importantbox}

\subsection{拓展阅读}

\begin{itemize}
    \item LeCun et al., ``Gradient-Based Learning Applied to Document Recognition'' (1998) —— LeNet-5 原始论文
    \item Krizhevsky et al., ``ImageNet Classification with Deep Convolutional Neural Networks'' (2012) —— AlexNet
    \item Dosovitskiy et al., ``An Image is Worth 16x16 Words'' (2021) —— Vision Transformer (ViT)
    \item CS231N 课程主页：\url{https://cs231n.stanford.edu/}
\end{itemize}

\end{document}
